\section{Additional Experiments}
\label{sec:appendix_experiments}

In this section, we provide additional experimental results supporting the main
claims of the paper. The experiments are designed to test four aspects of the
theory: (i) whether standard Langevin approximations reproduce discrete SGD,
(ii) whether the local covariance formula predicts the observed variance
profile in the Hessian eigenbasis, (iii) whether flat and sharp directions
exhibit distinct dynamical regimes, and (iv) whether the stochastic component
of SGD can be causally controlled. Unless stated otherwise, all experiments use
fixed random seeds, save the full configuration and raw intermediate artifacts,
and generate figures directly from saved outputs.

\subsection{Falsification of Langevin Approximations}

We first compare discrete SGD with standard and modified Langevin
approximations. The goal of these experiments is not merely to fit trajectories,
but to test whether the qualitative behavior of the stochastic process is
captured at finite learning rate. In the toy setting, the dynamics are chosen so
that the standard Langevin approximation predicts a stationary variance in a
regime where the exact discrete SGD recursion does not. This provides a direct
counterexample to the view that SGD can always be faithfully represented as
Brownian motion in a fixed potential.

\begin{figure}[h]
    \centering
    \includegraphics[width=0.7\linewidth]{images/new_exps_igor/exp1a_toy_full.png}
    \caption{Toy example showing qualitative discrepancies between SGD and Langevin dynamics.}
    \label{fig:appendix_exp1a_toy}
\end{figure}

We then repeat the comparison on the MLP-386 MNIST model. All methods are
initialized from the same reference point. For SGD, minibatches are sampled with
replacement. For the Langevin baselines, the deterministic drift is given by the
full-gradient direction, while the stochastic term is calibrated from minibatch
gradient fluctuations. We sweep the learning rate and compare the ensemble
variance, the mean trajectory, and the distance of the Langevin approximations
from the SGD ensemble. At small learning rates the methods remain close, while
finite learning rates produce systematic deviations from SGD. In the full run,
the modified Langevin correction improves the behavior only in some regimes and
does not uniformly recover the SGD dynamics, emphasizing that the discrete SGD
description is not simply reducible to a corrected Brownian model.

\begin{figure}[h]
    \centering
    \includegraphics[width=0.7\linewidth]{images/new_exps_igor/exp1b_mlp386_full.png}
    \caption{Comparison of SGD, standard Langevin, and modified Langevin dynamics on a small neural network. Standard Langevin deviates from SGD at finite learning rates, while the modified Langevin correction does not consistently recover SGD behavior in the full setup.}
    \label{fig:appendix_exp1b_mlp386}
\end{figure}

\subsection{Quantitative Validation of the Discrete Theory}

We next test the local covariance prediction derived from the discrete
Fokker--Planck analysis. Starting from a trained reference point, we compute the
full Hessian of the empirical loss and rotate trajectories into the eigenbasis
of the mean Hessian. The theoretical prediction is computed from independently
estimated gradient-noise statistics and Hessian eigenvalues; it is not obtained
by fitting the observed trajectory variance. We then compare the predicted and
measured variances across eigendirections.

The agreement is strongest in the leading eigendirections, where the local
quadratic approximation is most informative. In the full run, the predicted and
measured variance profiles have high log-correlation, and held-out checks show
that the agreement is not explained by a single post-hoc rescaling. We also
estimate the proportionality constant $\gamma$ across multiple sampler seeds and
find it to be stable. Finally, we directly check the Hessian-noise structure in
the mean-Hessian eigenbasis: the diagonal mass dominates the off-diagonal mass,
supporting the approximate diagonality assumption used in the derivation.

\begin{figure}[h]
    \centering
    \includegraphics[width=0.7\linewidth]{images/new_exps_igor/exp2_eq32_full.png}
    \caption{Quantitative validation of the theoretical predictions. Predicted and measured variances across eigendirections show strong structural agreement.}
    \label{fig:appendix_exp2_eq32}
\end{figure}

\subsection{Sampling Effects}

Our theoretical master equation assumes sampling with replacement, so that
minibatches are independent across steps. Standard epoch-based training often
uses sampling without replacement, which introduces temporal correlations
between updates. To assess the relevance of this distinction, we compare
trajectory ensembles generated from the same reference point and optimizer
hyperparameters under both sampling schemes. The experiment measures the
resulting covariance and trajectory displacement in the same projection space.
The results show that the sampling rule changes the stochastic dynamics,
confirming that replacement sampling is not merely an implementation detail but
part of the probabilistic model being tested.

\begin{figure}[h]
    \centering
    \includegraphics[width=0.7\linewidth]{images/new_exps_igor/exp3_sampling_full.png}
    \caption{Effect of sampling strategy on SGD dynamics.}
    \label{fig:appendix_exp3_sampling}
\end{figure}

\subsection{Scaling Diagnostics}

Exact dense Hessian computations are feasible only for small models. To probe
whether the sharp--flat decomposition persists beyond the smallest setting, we
include a lightweight scaling diagnostic based on Hessian-vector products and
top-curvature directions. This experiment does not aim to reproduce the full
dense-Hessian analysis at large scale. Instead, it checks whether a small number
of high-curvature directions can still be identified and whether projected
trajectory variance behaves differently along those directions compared with
random low-curvature directions.

\begin{figure}[h]
    \centering
    \includegraphics[width=0.7\linewidth]{images/new_exps_igor/exp4_scaling_full.png}
    \caption{Scaling diagnostics across configurations.}
    \label{fig:appendix_exp4_scaling}
\end{figure}

\subsection{Non-Stationarity and Learning-Rate Scaling}

The theory predicts that near non-stationary points the mean trajectory drifts
through a fluctuating local curvature field, producing additional sources of
spreading. We therefore compare SGD trajectories starting from an early or
mid-training reference point, where the full gradient is not zero. The measured
trajectory spread is compared with stationary and non-stationary approximations.
This experiment tests whether transient dynamics contain variance growth not
explained by a stationary local model.

\begin{figure}[h]
    \centering
    \includegraphics[width=0.7\linewidth]{images/new_exps_igor/exp5_nonstationary_full.png}
    \caption{Non-stationary behavior of SGD: certain directions exhibit persistent spreading over time.}
    \label{fig:appendix_exp5_nonstationary}
\end{figure}

We additionally sweep the learning rate and measure variance growth in two
Hessian-defined buckets: nearly-flat directions and sharp directions. The
important comparison is not whether sharp directions reach a mathematically
perfect plateau in finite time, but whether their late-time growth is much
weaker than the non-stationary spreading observed in flat or weakly-curved
directions. The full run shows learning-rate dependence of the variance-growth
statistics and effective confinement of the sharp directions, consistent with
the predicted separation between diffusive and rigid modes.

\begin{figure}[h]
    \centering
    \includegraphics[width=0.7\linewidth]{images/new_exps_igor/exp5_lr_scaling_full.png}
    \caption{Learning-rate scaling of SGD dynamics. Flat directions exhibit slow diffusion with learning-rate dependence, while sharp directions remain effectively confined.}
    \label{fig:appendix_exp5_lr_scaling}
\end{figure}

\subsection{Generalization and Geometry}

We also test whether displacement along Hessian-defined subspaces is directly
predictive of generalization. Starting from a common reference point, we run
multiple SGD trajectories and measure their final displacement in flat and sharp
eigenspaces, together with train loss, test loss, and generalization gap. This
experiment is exploratory: its purpose is to determine whether the local
diffusion picture immediately translates into a simple generalization predictor.
The observed correlations are weak, suggesting that the theory primarily
describes local stochastic dynamics rather than providing a standalone
generalization law.

\begin{figure}[h]
    \centering
    \includegraphics[width=0.7\linewidth]{images/new_exps_igor/exp6_generalization_full.png}
    \caption{Relationship between diffusion and generalization.}
    \label{fig:appendix_exp6_generalization}
\end{figure}

\subsection{Controlling Diffusion}

To test whether the theoretically identified directions are causally relevant,
we modify SGD updates by suppressing stochastic components in selected
subspaces. The selected subspaces include flat directions, sharp directions, and
directions predicted by the local theory to have high variance. Suppressing the
noise in these directions reduces the corresponding projected variance, showing
that the measured spread is not a plotting artifact but is directly controlled
by the minibatch-noise component of the update.

\begin{figure}[h]
    \centering
    \includegraphics[width=0.7\linewidth]{images/new_exps_igor/exp7_diffusion_control_full.png}
    \caption{Suppressing diffusion in selected eigendirections constrains the dynamics.}
    \label{fig:appendix_exp7_diffusion_control}
\end{figure}

Conversely, we amplify the stochastic component in selected directions. This
experiment checks the opposite intervention: if the selected subspace is
responsible for the observed diffusion, increasing the noise there should
increase the trajectory spread. The resulting curves confirm that the variance
can be increased in a targeted way, although this does not necessarily improve
optimization or generalization.

\begin{figure}[h]
    \centering
    \includegraphics[width=0.7\linewidth]{images/new_exps_igor/exp8_diffusion_amplification_full.png}
    \caption{Amplifying diffusion in selected directions enhances stochastic spreading.}
    \label{fig:appendix_exp8_diffusion_amplification}
\end{figure}

\subsection{Theory-Guided Denoising}

As a further intervention, we construct a theory-guided denoised update. At each
step, we decompose the minibatch gradient into the full-gradient component and
the stochastic residual, then remove the residual component in selected
eigendirections. This is not proposed as an efficient optimizer, since it uses
full-gradient information, but it provides a controlled test of the theory. The
experiment shows that removing noise in the predicted high-variance subspace
substantially reduces projected variance while leaving the deterministic drift
intact.

\begin{figure}[h]
    \centering
    \includegraphics[width=0.7\linewidth]{images/new_exps_igor/exp9_theory_denoised_sgd_full.png}
    \caption{Theory-guided denoising of SGD updates alters the dynamics in a predictable way.}
    \label{fig:appendix_exp9_theory_denoised}
\end{figure}

\subsection{Causal Control of Noise via Alpha Interpolation}

Finally, we introduce a direct interpolation between full-gradient descent and
SGD. For each minibatch we compute
\[
g_{\mathrm{batch}} = g_{\mathrm{full}} + \xi,
\]
where $\xi$ is the minibatch-noise component. We then update using
\[
g_\alpha = g_{\mathrm{full}} + \alpha \xi .
\]
Thus $\alpha=0$ corresponds to deterministic full-gradient descent from the
same reference point, while $\alpha=1$ recovers the SGD update. This experiment
provides a causal test of the role of minibatch noise: increasing $\alpha$
directly increases the ensemble variance of trajectories. In the raw-noise
setting, intermediate noise levels can match or improve the final test loss
relative to $\alpha=1$, indicating that the effect of stochasticity is
quantitative and regime-dependent rather than universally beneficial.

\begin{figure}[h]
    \centering
    \includegraphics[width=0.7\linewidth]{images/new_exps_igor/exp10_noise_alpha_sweep_full.png}
    \caption{Alpha sweep controlling the magnitude of the stochastic component in SGD. Increasing $\alpha$ increases trajectory spread and ensemble variance.}
    \label{fig:appendix_exp10_alpha}
\end{figure}

We also consider a normalized variant in which the noise vector is rescaled to
match the norm of the full gradient before applying the coefficient $\alpha$.
This isolates the effect of relative noise strength. The resulting trajectories
show a particularly clear monotonic increase in ensemble variance as $\alpha$
increases. In this setup, large normalized noise substantially increases
trajectory spread and can degrade test performance, demonstrating that the
stochastic component of SGD is a controllable driver of diffusion rather than an
external Brownian forcing independent of the loss landscape.

\begin{figure}[h]
    \centering
    \includegraphics[width=0.7\linewidth]{images/new_exps_igor/exp11_normalized_noise_alpha_sweep_full.png}
    \caption{Normalized alpha sweep. Controlled stochasticity directly modulates trajectory variance and affects generalization behavior.}
    \label{fig:appendix_exp11_alpha_normalized}
\end{figure}

